% Exemplo: quadro (bordas fechadas, conteúdo textual) com tabularx e colunas de largura fixa
% Compilar: latexmk -pdf quadro-requisitos.tex    Testado com: abntex2, TeX Live 2024 ou mais recente
\documentclass[12pt, a4paper, brazilian]{abntex2}
\usepackage[T1]{fontenc}
\usepackage{tabularx}                 % colunas X (largura automática)
\usepackage{array}                    % >{\raggedright\arraybackslash}

%%BEGIN:quadropreambulo
% Float "Quadro" (a abntex2 não o define): numeração e lista próprias (memoir)
\newfloat{quadro}{loq}{Quadro}
\newlistof{listofquadros}{loq}{Lista de quadros}
% linha da lista no formato da NBR 14724: "Quadro 1 – Título ..... 5"
\newlistentry{quadro}{loq}{0}
\renewcommand{\cftquadroname}{Quadro\space}
\renewcommand*{\cftquadroaftersnum}{\hfill\textendash\hfill}
\setlength{\cftquadronumwidth}{\cftfigurenumwidth}
\providecommand{\quadroautorefname}{Quadro}  % nome usado pelo \autoref
%%END:quadropreambulo

\begin{document}

Os requisitos estão no \autoref{qua:requisitos}.

%%BEGIN:quadrorequisitos
\begin{quadro}[htbp]
  \caption{Requisitos funcionais do sistema (exemplo)}\label{qua:requisitos}
  \renewcommand{\arraystretch}{1.25}
  \begin{tabularx}{\textwidth}{|c|>{\raggedright\arraybackslash}p{3.6cm}|X|c|}
    \hline
    \textbf{ID} & \textbf{Requisito}
      & \textbf{Descrição}
      & \textbf{Prioridade} \\
    \hline
    RF01 & Cadastrar usuário
      & O sistema deve permitir que o visitante crie uma conta informando
        nome, \emph{e-mail} e senha.
      & Alta \\
    \hline
    RF02 & Autenticar usuário
      & O sistema deve permitir o acesso por \emph{e-mail} e senha e
        bloquear a conta após cinco tentativas inválidas.
      & Alta \\
    \hline
    RF03 & Gerar relatório
      & O sistema deve gerar relatório mensal de atividades, exportável em
        PDF e em CSV.
      & Média \\
    \hline
  \end{tabularx}
  \fonte{Elaborado pelo próprio autor.}
\end{quadro}
%%END:quadrorequisitos

\end{document}
